\documentclass[11pt,a4paper]{article}
\usepackage[margin=2.2cm]{geometry}
\usepackage{booktabs,tabularx,xcolor,hyperref,enumitem}
\hypersetup{colorlinks,linkcolor=blue!50!black,urlcolor=blue!50!black}
\setlist{itemsep=1pt,topsep=2pt}
\newcommand{\code}[1]{\texttt{\small #1}}
\definecolor{ph1}{HTML}{4E79A7}\definecolor{ph2}{HTML}{F28E2B}\definecolor{ph3}{HTML}{59A14F}
\newcommand{\step}[5]{%
\par\medskip\noindent\textbf{#1}\par
\noindent\begin{tabularx}{\linewidth}{@{}lX@{}}
Do & #2\\
Why & #3\\
You get & #4\\
Where & #5\\
\end{tabularx}}

\title{How the HY Analyst works, step by step}
\author{FundamentalQuantmental}
\date{22 October 2026}

\begin{document}
\maketitle

\noindent\fcolorbox{gray}{gray!6}{\parbox{\dimexpr\linewidth-2\fboxsep-2\fboxrule}{%
\textbf{The whole method in one line.} Measure the issuer today, project it under each macro story, price what the bond should pay, and buy only where the market is wrong by more than the cost.

\smallskip
\textbf{Three phases, eleven steps.}
\textcolor{ph1}{\textbf{Measure}} (steps 1 to 4) $\rightarrow$
\textcolor{ph2}{\textbf{Forecast}} (steps 5 to 7) $\rightarrow$
\textcolor{ph3}{\textbf{Decide}} (steps 8 to 11).

\smallskip
\textbf{One rule over all of it.} Every number carries the date the market could first read it. No step reads a number from the future.}}

\section*{\textcolor{ph1}{Phase 1: Measure the issuer today}}

\step{1. Load the data with its release time}
{Pull financials (Compustat point-in-time, SEC filings), bond prices (TRACE), and the Treasury curve. Stamp each number with \code{available\_at}.}
{A backtest that reads a number before its release finds profit that never existed.}
{Quarterly records, every version of every quarter.}
{Methodology \S3; \code{hy/inputs.py}, \code{hy/sources/}}

\step{2. Frame the business}
{Pick one archetype (for example lease-heavy, cyclical, asset-heavy contracted) and up to two modifiers (sponsor, decline, concentration).}
{The archetype sets which EBITDA to use, how the business fails, and which events to watch.}
{\code{archetype}, \code{modifiers}, a one-line reason.}
{Framework F2}

\step{3. Measure health now}
{Adjust EBITDA for one-off items. Compute free cash flow, liquidity, leverage, coverage, covenant headroom, risk flags, and the rating gap.}
{These are the inputs every later step uses.}
{One state vector per issuer per filing.}
{Methodology \S4; \code{hy/state.py}}

\step{4. Check the two clocks}
{Put fixed charges (rent, debt, interest) and signed revenue on one yearly grid.}
{When the bills outlast the contracts, the bond is a bet on renewal.}
{Breakeven margin, renewal-dependent charges, largest shortfall.}
{Methodology \S4.7; \code{hy/clocks.py}}

\section*{\textcolor{ph2}{Phase 2: Forecast under each macro story}}

\step{5. Turn forecasts into macro paths}
{Read GDP, inflation, rate, and spread forecasts with their vintages. Build five paths: soft landing, recession, higher for longer, disinflation rally, stagflation. Weight each path by the forecasts, then update with the cycle phase.}
{The forecasters' own uncertainty sets the odds, not the analyst's mood.}
{Five paths of eight quarters, each with a probability; prior, forecast mass, and posterior recorded.}
{Methodology \S5.4; \code{hy/forecasts.py}, \code{hy/macro.py}}

\step{6. Run every case under every path}
{Build six cases (management, consensus, base, moderate downside, severe downside, upside). Run each one under each path. Refinance each maturity at the forward Treasury plus fair spread, or record a failure.}
{One case on one path hides the tail; the grid shows it.}
{30 projected paths of revenue, EBITDA, debt, cash, and headroom.}
{Methodology \S5.1; framework F6; \code{hy/forward.py}}

\step{7. Weight the cases and average}
{Weight the cases from the F6 table (later the fitted logistic). If the downside weight passes 0.5, shift every weight down one step (F6.6). Average across cases and paths.}
{The expected number must include how likely each story is.}
{Expected leverage and its spread, breach, liquidity, and refinance odds, per path and weighted.}
{Methodology \S5.3; framework F6.6}

\section*{\textcolor{ph3}{Phase 3: Decide what to own}}

\step{8. Price what the bond should pay}
{Spread from expected loss (default probability times loss), times the market price of risk, plus a model for the rest. Compare with the market spread. Flag cheap only outside the error band.}
{A bond is cheap against everything the model knows, not against a noisy point.}
{Fair spread, residual, cheap or rich flag.}
{Methodology \S6, \S7; \code{hy/recovery.py}, \code{hy/fairspread.py}}

\step{9. Convert to expected alpha}
{Per path: carry less the hedge, plus the part of the gap that closes, less expected default loss. Weight by path, add event P\&L, subtract trading cost.}
{The number you act on is return after risk and cost, not spread.}
{\code{exp\_alpha}, its parts, its spread across paths, and the information ratio.}
{Methodology \S7.5; \code{hy/alpha.py}}

\step{10. Signal, gate, and rank}
{Score nine signals. Pass the gate: a gap larger than its noise, a dated catalyst, alpha above cost, one observable that would prove the call wrong. Rank the whole universe and name the position each new one replaces.}
{A credit opinion is not a trade until the price is wrong by enough.}
{The force rank, with \code{replaces\_position} and \code{invalidation} for each entry.}
{Methodology \S8; framework F9, F10; \code{hy/signals.py}}

\step{11. Size for events, then check the record}
{For each dated event, list outcomes and their odds under each path; compute P\&L per outcome; cap size to the event budget. After each event and each exit, record what happened against what was predicted.}
{The model that calls a bond cheap must also say how much it loses on a miss, and must be scored.}
{Event matrix, limits, hedges, postmortem.}
{Methodology \S9 to \S11; framework F11, F13; \code{hy/events.py}}

\section*{What you supply, and what the system computes}
\begin{itemize}
\item \textbf{You:} the archetype, analyst cases and gaps to consensus, ratings, document terms, and any override in \code{data/analyst\_inputs/}.
\item \textbf{The system:} everything else, each number with its source and release time.
\item \textbf{Every hand-set number} has a decision page in \code{docs/decisions/} with the test that would change it.
\end{itemize}

\end{document}
